\chapter{ Reglas de inferencia }

  \section{Modos Clásicos}

  Podemos simplificar la notación asignando nombres a ciertas
  proposiciones; estos son conocidos como los 4 modos clásicos.

  \begin{definition}[title={Modus Ponens}]
  La restricción se propaga. Estructura:
  \begin{align}
  &p \to q \\
  &p \\
  \hline
  &\therefore q
  \end{align}
  \end{definition}

  \begin{definition}[title={Modus Tollens}]
  Ya que la restricción es $p \land \sim q$ y $\sim q = V$, no puede
  ocurrir que $p = V$, de modo que solo queda $\sim p$. Estructura:
  \begin{align}
  &p \to q \\
  &\sim q \\
  \hline
  &\therefore \sim p
  \end{align}
  \end{definition}

  \begin{definition}[title={Silogismo Hipotético}]
  Es básicamente una composición de restricciones. Estructura:
  \begin{align}
  &p \to q \\
  &q \to r \\
  \hline
  &\therefore p \to r
  \end{align}
  \end{definition}

  \begin{definition}[title={Silogismo Disyuntivo}]
  Un \textit{o} excluyente. Estructura:
  \begin{align}
  &p \lor q \\
  &\sim p \\
  \hline
  &\therefore q
  \end{align}
  \end{definition}

  \section{Argumento Válido}

  \begin{definition}[title={Argumento}]
  Un argumento tiene la forma $(P_1, P_2, \ldots) \therefore q$, o
  formalmente $(P_1 \land P_2 \land \cdots) \to q$, y diremos que es
  válido si no existe una valuación donde las premisas son verdaderas y
  la conclusión es falsa.
  \end{definition}

  Y aquí entra una aclaración muy importante: toda regla de inferencia se
  puede justificar mostrando que determinada fórmula es una tautología.
  Para ello, tome el \textit{modus ponens}:

  Para verificarlo construimos la fórmula $((p \to q) \land p) \to q$.
  Si esta expresión es una tautología, entonces la regla es válida.

  \begin{tabular}{c|c|c|c|c}
  $p$ & $q$ & $p \to q$ & $(p \to q) \land p$ & $((p \to q) \land p) \to q$ \\
  \hline
  V & V & V & V & V \\
  V & F & F & F & V \\
  F & V & V & F & V \\
  F & F & V & F & V \\
  \end{tabular}

  Esto a su vez es aplicar la tautología como patrón, de forma que su
  notación nos permite reemplazar bloques completos con la hipótesis ya
  probada. De este modo, si tenemos
  $((p \to q) \land (q \to r)) \land p) \to r$, entonces podemos
  reemplazar todo el primer bloque por $r$ sin perder verdad.

  \begin{definition}[title={Verdad bajo valuación}]
  Decimos que una fórmula es verdad bajo una valuación $v$ si al
  evaluarla con las reglas semánticas de los conectores el resultado es $V$.
  \end{definition}

  \begin{definition}[title={Argumento válido}]
  Un argumento es válido si no existe una valuación donde todas las
  premisas sean verdaderas y las conclusiones sean falsas.
  \end{definition}